\scanpage{195}
\mspage{II.43}
% archive sheet number: 194
% block: top, continuation of 6.13 | text
naissance à des groupes proanalytiques et proalgébriques
% block: top, completions | display
\[
\unsure{\pi_1^{\mathrm{arith}}|\Pi_1^{\mathrm{arith}}}(X,\unsure{\xi|\zeta})^{\mathrm{an}},\quad
\pi_1^{\mathrm{arith}}(X,\xi)^{\mathrm{alg}}
\]
% block: upper third, Lie algebras | text
et à des proalgèbres de Lie correspondantes.

% block: upper third, 6.14 | text
6.14. Supposons $X$ intègre et normal, alors on peut supposer les
$X_\lambda$ intègres, normaux (quitte à \deleted{les} remplacer $X_\lambda$
par le normalisé de l'image schématique fermée de $X$ dans $X_\lambda$
$\ldots$). \deleted{\unsure{Les}}
Comme \unsure{pour} tt morphisme dominant de t.f. de schémas normaux,
$X_\mu\mathbin{-}X_\lambda$ induit un \deleted{quasi} hom.
\uline{quasi-surjectif}
$\pi_1(X_\mu,\xi_\mu)\mathbin{-}\pi_1(X_\lambda,\xi_\lambda)$,
i.e. on est dans le cas de 6.12. En particulier, les \deleted{syst.} hom.
de transition dans \deleted{$\unsure{\Pi}$}
$G=\pi_1^{\mathrm{arith}}(X,\xi)^{\mathrm{an}}$ et
$\widetilde{G}=\pi_1^{\mathrm{arith}}(X,\xi)^{\mathrm{alg}}$ sont des
quasi-épimorphismes, et les hom. de transition dans $\mathfrak{g}$ et
$\widetilde{\mathfrak{g}}$ sont des épimorphismes (i.e.
$\mathfrak{g},\unsure{\widehat{\mathfrak{g}}|\widetilde{\mathfrak{g}}}$
sont des pro-objets \uline{stricts}).

% block: lower third, 6.15 and interlinear corrections | text
6.15. Notons que \deleted{si $\ell$ est inversible sur $X$.}
\inserted{si $\ell$ est inversible sur $X$ (\unsure{seul cas intéressant}
$\ldots$) alors}
$\Pi\ (=\pi_1^{\mathrm{arith}}(X,\xi))$
n'est \deleted{\unsure{jamais}} ``discret'', i.e.
\deleted{\unsure{il y a d}} $\dim\mathfrak{g}>0$,
\deleted{$\unsure{\dim\mathfrak{g}}$} ou encore
$\dim\widetilde{\mathfrak{g}}>0$. (En fait, on verra que ces dimensions
sont \uline{tjs} infinies).

% block: foot, argument beginning | text
En effet, il suffit de prouver que si \deleted{et $\ell$ y est inversible}
$X$ est de type fini sur $\operatorname{Spec}\mathbb{Z}$, et muni d'un
pt géom. $\xi$, \inserted{et si $\ell$ y est inversible,} alors,
\deleted{$\unsure{\Pi_1(X,|\pi_1(X,}$} \unsure{posons}
$\pi=\pi_1(X,\xi)$, $\pi^{\mathrm{alg}}$ n'est jamais discret. Or
